\section{Дополнительный диагностический протокол}
Следующие проверки были заданы после получения основных результатов. Исходные правила остаются зафиксированными; сохранены все 48 выборок и все заданные диагностические условия. Интервалы описательные, парные 95\%, по независимым обучающим выборкам.

\begin{table}[p]\caption{Средняя стоимость в тысячах для каждого механизма и диагностического сравнения.}\centering\small\setlength{\tabcolsep}{3pt}\begin{tabular}{p{.40\textwidth}rrrr}\toprule
Метод & Пост. & Длительность & Возвраты & Совместное\\\midrule
Без действия & 23.6 & 211.7 & 147.8 & 375.2\\
Статическое планирование & 18.0 & 111.0 & 92.8 & 309.2\\
Отбор по прогнозу & 18.0 & 111.0 & 88.8 & 292.5\\
Правило по силе зависимости & 18.9 & 111.0 & 73.0 & 309.2\\
Сохранение нагрузочных компонент & 19.9 & 111.0 & 57.4 & 288.3\\
Ресурс при 20 & 17.7 & 111.0 & 92.6 & 309.2\\
Возвраты при 20 & 22.6 & 202.7 & 45.3 & 284.1\\
Известный генератор & 17.8 & 111.0 & 45.3 & 284.1\\
Ретроспективная граница & 17.7 & 111.0 & 45.3 & 284.1\\
\bottomrule\end{tabular}\end{table}
\begin{table}[p]\caption{Превышение стоимости относительно планировщика с известным генератором: среднее и 95\% интервал, в тысячах. Отрицательные значения возможны из-за конечной выборки и иных сценариев планирования.}\centering\small\setlength{\tabcolsep}{3pt}\begin{tabular}{llr}\toprule
Механизм & Метод & Превышение [интервал]\\\midrule
Постоянное & Статическое планирование & 0.17 [-0.08, 0.42]\\
Постоянное & Отбор по прогнозу & 0.17 [-0.08, 0.42]\\
Постоянное & Правило по силе зависимости & 1.09 [-0.19, 2.36]\\
Постоянное & Сохранение нагрузочных компонент & 2.13 [0.41, 3.85]\\
Длительность & Статическое планирование & 0.00 [0.00, 0.00]\\
Длительность & Отбор по прогнозу & 0.00 [0.00, 0.00]\\
Длительность & Правило по силе зависимости & 0.00 [0.00, 0.00]\\
Длительность & Сохранение нагрузочных компонент & 0.00 [0.00, 0.00]\\
Возвраты & Статическое планирование & 47.52 [45.70, 49.34]\\
Возвраты & Отбор по прогнозу & 43.52 [34.63, 52.42]\\
Возвраты & Правило по силе зависимости & 27.66 [12.12, 43.20]\\
Возвраты & Сохранение нагрузочных компонент & 12.04 [-1.87, 25.94]\\
Совместное & Статическое планирование & 25.05 [23.43, 26.68]\\
Совместное & Отбор по прогнозу & 8.32 [0.44, 16.19]\\
Совместное & Правило по силе зависимости & 25.05 [23.43, 26.68]\\
Совместное & Сохранение нагрузочных компонент & 4.20 [-2.05, 10.45]\\
\bottomrule\end{tabular}\end{table}
\begin{table}[p]\caption{Эффект вмешательства: решения зафиксированы при предполагаемом сокращении 50\%; показаны фактические сокращения. Средние стоимости в тысячах.}\centering\small\setlength{\tabcolsep}{3pt}\begin{tabular}{llrrrrr}\toprule
Механизм & Эффект & Стат. & Отбор & Нагрузка & Ресурс & Возвраты\\\midrule
Постоянное & 30\% & 18.0 & 18.0 & 20.1 & 17.7 & 23.2\\
Постоянное & 50\% & 18.0 & 18.0 & 19.9 & 17.7 & 22.6\\
Постоянное & 70\% & 18.0 & 18.0 & 19.7 & 17.7 & 22.0\\
Длительность & 30\% & 111.0 & 111.0 & 111.0 & 111.0 & 206.5\\
Длительность & 50\% & 111.0 & 111.0 & 111.0 & 111.0 & 202.7\\
Длительность & 70\% & 111.0 & 111.0 & 111.0 & 111.0 & 199.1\\
Возвраты & 30\% & 92.8 & 91.1 & 77.1 & 92.6 & 71.7\\
Возвраты & 50\% & 92.8 & 88.8 & 57.4 & 92.6 & 45.3\\
Возвраты & 70\% & 92.8 & 87.8 & 47.6 & 92.6 & 32.2\\
Совместное & 30\% & 309.2 & 316.8 & 319.2 & 309.2 & 321.1\\
Совместное & 50\% & 309.2 & 292.5 & 288.3 & 309.2 & 284.1\\
Совместное & 70\% & 309.2 & 268.1 & 258.0 & 309.2 & 248.1\\
\bottomrule\end{tabular}\end{table}
\begin{table}[p]\caption{Чувствительность к смеси: средняя нормированная разность нагрузочного метода и сравнения; отрицательное значение благоприятно. Меняется вес постоянного механизма, остальные три делят оставшийся вес поровну.}\centering\small\setlength{\tabcolsep}{3pt}\begin{tabular}{p{.37\textwidth}rrrr}\toprule
Сравнение & 0 & 0.25 & 0.5 & 1\\\midrule
Статическое планирование & -0.099 & -0.054 & -0.008 & 0.083\\
Отбор по прогнозу & -0.075 & -0.036 & 0.004 & 0.083\\
Ресурс при 20 & -0.099 & -0.050 & -0.002 & 0.095\\
Возвраты при 20 & -0.114 & -0.114 & -0.113 & -0.113\\
Правило по силе зависимости & -0.054 & -0.030 & -0.005 & 0.044\\
Известный генератор & 0.030 & 0.045 & 0.060 & 0.090\\
\bottomrule\end{tabular}\end{table}
\begin{table}[p]\caption{Чувствительность критерия события для возвратов: процент траекторий, достигших критерия; исходные зафиксированные правила.}\centering\small\setlength{\tabcolsep}{3pt}\begin{tabular}{rrrrrr}\toprule
WIP & Длительность & Без действия & Стат. & Нагрузка & Возвраты\\\midrule
80 & 100 & 61.6 & 50.5 & 39.9 & 37.5\\
80 & 200 & 59.0 & 43.6 & 37.3 & 34.2\\
80 & 300 & 55.6 & 38.2 & 34.5 & 33.3\\
100 & 100 & 54.2 & 39.2 & 35.6 & 33.7\\
100 & 200 & 50.2 & 35.8 & 34.0 & 33.3\\
100 & 300 & 46.9 & 34.4 & 33.3 & 32.5\\
120 & 100 & 43.9 & 34.4 & 34.0 & 33.3\\
120 & 200 & 42.7 & 34.0 & 33.2 & 32.1\\
120 & 300 & 41.3 & 33.9 & 30.4 & 29.0\\
\bottomrule\end{tabular}\end{table}
\begin{table}[p]\caption{Чувствительность стоимости для возвратов: горизонт и вес конечного остатка. Средние стоимости в тысячах; реакции повторно не выбираются.}\centering\small\setlength{\tabcolsep}{3pt}\begin{tabular}{rrrrrr}\toprule
Горизонт & Вес & Стат. & Отбор & Нагрузка & Возвраты\\\midrule
1400 & 0 & 77.0 & 74.2 & 52.2 & 43.7\\
1400 & 100 & 81.9 & 78.7 & 53.8 & 44.3\\
1400 & 200 & 86.8 & 83.2 & 55.4 & 44.9\\
1600 & 0 & 85.9 & 82.4 & 55.1 & 44.6\\
1600 & 100 & 89.7 & 85.9 & 56.3 & 44.9\\
1600 & 200 & 93.5 & 89.5 & 57.4 & 45.3\\
1800 & 0 & 92.8 & 88.8 & 57.4 & 45.3\\
1800 & 100 & 95.8 & 91.6 & 58.4 & 45.6\\
1800 & 200 & 98.7 & 94.3 & 59.4 & 45.9\\
\bottomrule\end{tabular}\end{table}
Архив содержит полный выбор действий по каждой выборке, распределения порогов, исходные условия эффективности и чувствительности, выявленные сети, проверки компилятора и машиночитаемую карточку реакции. 4032 новых запусков планирования с известным генератором; 9156 новых запусков с изменённым эффектом и 4668 точно повторно использованных траекторий без активации. Повторное использование допустимо, поскольку эффект меняет только события после активации, которая на этих траекториях не происходит.
\clearpage
\begin{table}[p]\caption{Числа включений и предотвращений при зафиксированных нагрузочных реакциях. Последний столбец считает включения на траекториях без тяжёлого события в сопоставленном запуске без действия; это отличается от включения в спокойном контроле.}\centering\small\setlength{\tabcolsep}{3pt}\begin{tabular}{lrrrr}\toprule
Механизм & Риск & Ранние & Предотвращено & \shortstack{Включение без\\события в базе}\\\midrule
Постоянное & 117 & 109 & 74 & 239\\
Длительность & 382 & 382 & 138 & 29\\
Возвраты & 289 & 289 & 93 & 95\\
Совместное & 421 & 421 & 12 & 18\\
\bottomrule\end{tabular}\end{table}\clearpage
